\documentclass[border=10pt,tikz]{standalone}
\usepackage{tikz}
\usepackage{fontspec}
\setmainfont{Apple SD Gothic Neo}
\setsansfont{Apple SD Gothic Neo}
\usetikzlibrary{positioning,arrows.meta,shapes,fit,backgrounds,calc}
\begin{document}
% _design.tex — shared TikZ design system for all blog diagrams
% Inputs nothing on its own — meant to be \input{} from each diagram .tex
% AFTER \begin{document}, BEFORE \begin{tikzpicture}.
%
% Provides a unified visual language across GoF / DSA / EMC++ / Beautiful series:
%   - Colors: muted, accessible, color-coded by role
%   - Node styles: consistent sizing, rounded, subtle borders
%   - Edge styles: UML-style inheritance, aggregation, plus dashed alternates
%
% Usage:
%   \documentclass[border=10pt,tikz]{standalone}
%   \usepackage{tikz}
%   \usetikzlibrary{positioning,arrows.meta,shapes,fit,backgrounds}
%   \begin{document}
%   \input{../_design.tex}     % adjust relative path
%   \begin{tikzpicture}[blog]  % apply 'blog' style
%   ...

\definecolor{absfill}{RGB}{220,232,247}   % cool blue
\definecolor{absbord}{RGB}{72,118,182}
\definecolor{confill}{RGB}{249,222,222}   % warm red
\definecolor{conbord}{RGB}{197,93,93}
\definecolor{impfill}{RGB}{223,238,222}   % soft green
\definecolor{impbord}{RGB}{99,154,93}
\definecolor{clifill}{RGB}{251,239,205}   % pale yellow
\definecolor{clibord}{RGB}{200,167,72}
\definecolor{hifill} {RGB}{253,224,200}   % accent orange
\definecolor{hibord} {RGB}{210,135,68}
\definecolor{nodefill}{RGB}{250,250,250}
\definecolor{nodebord}{RGB}{120,120,120}
\definecolor{edgec}  {RGB}{80,80,80}
\definecolor{edgesc} {RGB}{145,145,145}
\definecolor{groupbord}{RGB}{200,200,200}

\tikzset{
  blog/.style={
    font=\sffamily\small,
    every node/.style={font=\sffamily\small}
  },
  % --- node roles ---
  cls/.style={
    rectangle, rounded corners=2pt,
    draw=nodebord, line width=0.4pt,
    minimum height=8mm, minimum width=28mm,
    inner sep=3pt, fill=nodefill, align=center
  },
  abs/.style={cls, draw=absbord, fill=absfill, font=\sffamily\itshape\small},
  con/.style={cls, draw=conbord, fill=confill},
  imp/.style={cls, draw=impbord, fill=impfill},
  cli/.style={cls, draw=clibord, fill=clifill},
  hi/.style={cls, draw=hibord, fill=hifill, font=\sffamily\bfseries\small},
  % --- circle (for trees / graph nodes) ---
  vx/.style={
    circle, draw=absbord, line width=0.4pt,
    minimum size=8mm, fill=absfill, inner sep=0pt
  },
  vxc/.style={vx, draw=conbord, fill=confill},
  vxh/.style={vx, draw=hibord, fill=hifill, font=\sffamily\bfseries},
  % --- decision diamond ---
  q/.style={
    diamond, draw=clibord, line width=0.4pt, fill=clifill,
    inner sep=2pt, aspect=2.4, align=center
  },
  % --- edges ---
  inh/.style={                                  % UML inheritance
    -{Triangle[length=3mm,width=3mm,open]},
    draw=edgec, line width=0.5pt
  },
  agg/.style={                                  % UML aggregation
    {Diamond[length=2.5mm,width=2.5mm,open]}-{Stealth[length=2.2mm]},
    draw=edgec, line width=0.45pt
  },
  uses/.style={                                 % directed reference
    -{Stealth[length=2.2mm]}, draw=edgec, line width=0.5pt
  },
  arr/.style={                                  % plain arrow (graphs/trees)
    -{Stealth[length=2mm]}, draw=edgec, line width=0.45pt
  },
  edge/.style={                                 % undirected (trees)
    draw=edgec, line width=0.45pt
  },
  alt/.style={                                  % dashed: similar/alt
    -{Stealth[length=2mm]}, dashed, draw=edgesc, line width=0.45pt
  },
  ret/.style={                                  % dashed creates/returns
    -{Stealth[length=2mm]}, dashed, draw=edgesc, line width=0.45pt
  },
  thr/.style={                                  % threading link (red dashed)
    ->, dashed, draw=conbord, line width=0.5pt
  },
  % --- group / region ---
  group/.style={
    rectangle, rounded corners=4pt,
    draw=groupbord, line width=0.4pt,
    inner sep=10pt, fill=black!2
  },
  glabel/.style={font=\sffamily\bfseries\small, anchor=north west, inner sep=2pt}
}

% Korean font convenience macro — included by every Hangul diagram
\providecommand{\KOR}{}

\begin{tikzpicture}[blog,
  box/.style={cls, minimum width=66mm, text width=62mm, minimum height=15mm, align=center, font=\sffamily\footnotesize, inner sep=4pt},
  uclassB/.style={box, fill=absfill, draw=absbord},
  driverB/.style={box, fill=impfill, draw=impbord},
  devB/.style={box, fill=hifill, draw=hibord, minimum width=50mm, text width=46mm},
  stg/.style={cls, minimum width=46mm, text width=42mm, minimum height=14mm, align=center, font=\sffamily\footnotesize, inner sep=4pt},
  rel/.style={font=\sffamily\scriptsize, text=mutec, fill=white, inner sep=1.5pt},
  hdr/.style={font=\sffamily\bfseries\small, text=inkc, anchor=west}
]

% ── Part 1: three elements ──
\node[hdr] at (-3.0, 9.4) {세 요소의 관계 (같은 MMC driver를 쓰는 controller 두 개)};

\node[uclassB] (uc) at (4.5, 7.6) {%
  \textbf{uclass} — 같은 종류 device의 공통 interface\\[2pt]
  \scriptsize 예: \texttt{UCLASS\_MMC}, \texttt{UCLASS\_GPIO}};

\node[devB] (d0) at (0.5, 4.6) {%
  \textbf{device} (\texttt{udevice})\\[2pt]
  MMC controller 0의 instance};
\node[devB] (d1) at (8.5, 4.6) {%
  \textbf{device} (\texttt{udevice})\\[2pt]
  MMC controller 1의 instance};

\node[driverB] (drv) at (4.5, 1.6) {%
  \textbf{driver} — hardware-specific 동작 구현\\[2pt]
  \scriptsize 하나의 구현을 여러 instance가 사용};

\draw[uses] (d0) -- node[rel] {소속} (uc);
\draw[uses] (d1) -- node[rel] {소속} (uc);
\draw[uses] (d0) -- node[rel] {사용} (drv);
\draw[uses] (d1) -- node[rel] {사용} (drv);

% ── Part 2: bind and probe ──
\node[hdr] at (-3.0, -0.6) {DT 구성에서 device가 쓸 수 있게 되기까지};

\node[stg, draw=prpbord, fill=prpfill] (dt) at (-0.5, -2.4) {%
  \textbf{DT node}\\[2pt]
  \texttt{compatible} $\leftrightarrow$\\[2pt] driver의 matching data};
\node[stg, draw=clibord, fill=clifill] (bind) at (4.5, -2.4) {%
  \textbf{bind}\\[2pt]
  instance와 부모-자식 관계를\\[2pt] driver model에 등록};
\node[stg, draw=impbord, fill=impfill] (probe) at (9.5, -2.4) {%
  \textbf{probe}\\[2pt]
  driver와 관련 자원을 초기화해\\[2pt] 사용 가능 상태로};

\draw[flow] (dt) -- (bind);
\draw[flow] (bind) -- (probe);

\node[caption, text width=150mm] at (4.5, -4.0) {probe 시점은 uclass lookup, \texttt{DM\_FLAG\_PROBE\_AFTER\_BIND}, 단계별 초기화 정책에 따라 달라집니다.\\[2pt] bind가 됐다고 해서 clock·reset·pinctrl·parent bus까지 정상이라는 뜻은 아닙니다.};

\end{tikzpicture}
\end{document}
